\documentclass[10pt,a4paper]{amsart}
\usepackage{amsmath,amssymb,amsthm,graphicx}
\usepackage{verbatim}
\newenvironment{reference}{\comment}{\endcomment}
%\newenvironment{reference}{}{}
\newenvironment{slogan}{\comment}{\endcomment}
\newenvironment{history}{\comment}{\endcomment}

% For commutative diagrams we use Xy-pic
\usepackage[all]{xy}

% We use 2cell for 2-commutative diagrams.
\xyoption{2cell}
% Theorem environments.
%
\theoremstyle{plain}
\newtheorem{theorem}[subsection]{Theorem}
\newtheorem{proposition}[subsection]{Proposition}
\newtheorem{lemma}[subsection]{Lemma}

\theoremstyle{definition}
\newtheorem{definition}[subsection]{Definition}
\newtheorem{example}[subsection]{Example}
\newtheorem{exercise}[subsection]{Exercise}
\newtheorem{situation}[subsection]{Situation}

\theoremstyle{remark}
\newtheorem{remark}[subsection]{Remark}
\newtheorem{remarks}[subsection]{Remarks}

\numberwithin{equation}{subsection}

% Macros
%
\def\lim{\mathop{\mathrm{lim}}\nolimits}
\def\colim{\mathop{\mathrm{colim}}\nolimits}
\def\Spec{\mathop{\mathrm{Spec}}}
\def\Hom{\mathop{\mathrm{Hom}}\nolimits}
\def\Ext{\mathop{\mathrm{Ext}}\nolimits}
\def\SheafHom{\mathop{\mathcal{H}\!\mathit{om}}\nolimits}
\def\SheafExt{\mathop{\mathcal{E}\!\mathit{xt}}\nolimits}
\def\Sch{\mathit{Sch}}
\def\Mor{\mathop{\mathrm{Mor}}\nolimits}
\def\Ob{\mathop{\mathrm{Ob}}\nolimits}
\def\Sh{\mathop{\mathit{Sh}}\nolimits}
\def\NL{\mathop{N\!L}\nolimits}
\def\CH{\mathop{\mathrm{CH}}\nolimits}
\def\proetale{{pro\text{-}\acute{e}tale}}
\def\etale{{\acute{e}tale}}
\def\QCoh{\mathit{QCoh}}
\def\Ker{\mathop{\mathrm{Ker}}}
\def\Im{\mathop{\mathrm{Im}}}
\def\Coker{\mathop{\mathrm{Coker}}}
\def\Coim{\mathop{\mathrm{Coim}}}

% Boxtimes
%
\DeclareMathSymbol{\boxtimes}{\mathbin}{AMSa}{"02}

%
% Macros for moduli stacks/spaces
%
\def\QCohstack{\mathcal{QC}\!\mathit{oh}}
\def\Cohstack{\mathcal{C}\!\mathit{oh}}
\def\Spacesstack{\mathcal{S}\!\mathit{paces}}
\def\Quotfunctor{\mathrm{Quot}}
\def\Hilbfunctor{\mathrm{Hilb}}
\def\Curvesstack{\mathcal{C}\!\mathit{urves}}
\def\Polarizedstack{\mathcal{P}\!\mathit{olarized}}
\def\Complexesstack{\mathcal{C}\!\mathit{omplexes}}
% \Pic is the operator that assigns to X its picard group, usage \Pic(X)
% \Picardstack_{X/B} denotes the Picard stack of X over B
% \Picardfunctor_{X/B} denotes the Picard functor of X over B
\def\Pic{\mathop{\mathrm{Pic}}\nolimits}
\def\Picardstack{\mathcal{P}\!\mathit{ic}}
\def\Picardfunctor{\mathrm{Pic}}
\def\Deformationcategory{\mathcal{D}\!\mathit{ef}}
\usepackage{cleveref}
\crefname{originalreference}{original reference}{original references}
\Crefname{originalreference}{Original reference}{Original references}
\makeatletter
\newlabel{definition-groupoid-in-functors}{{06K3}{source}{}{sourcepointer-1}{}}
\newlabel{definition-groupoid-in-functors@cref}{{[originalreference][0][]06K3}{[1][0][]source}{}{}{}}
\makeatother
\crefname{originalreference}{original reference}{original references}
\Crefname{originalreference}{Original reference}{Original references}
\makeatletter
\newlabel{lemma-characterize-minimal-groupoid-in-functors}{{06KN}{source}{}{sourcepointer-2}{}}
\newlabel{lemma-characterize-minimal-groupoid-in-functors@cref}{{[originalreference][0][]06KN}{[1][0][]source}{}{}{}}
\makeatother
\crefname{originalreference}{original reference}{original references}
\Crefname{originalreference}{Original reference}{Original references}
\makeatletter
\newlabel{lemma-deformation-functor-diagonal}{{06L6}{source}{}{sourcepointer-3}{}}
\newlabel{lemma-deformation-functor-diagonal@cref}{{[originalreference][0][]06L6}{[1][0][]source}{}{}{}}
\makeatother
\crefname{originalreference}{original reference}{original references}
\Crefname{originalreference}{Original reference}{Original references}
\makeatletter
\newlabel{lemma-surjective-morphism-prorepresentable-functor}{{06KP}{source}{}{sourcepointer-4}{}}
\newlabel{lemma-surjective-morphism-prorepresentable-functor@cref}{{[originalreference][0][]06KP}{[1][0][]source}{}{}{}}
\makeatother
\title{Self-equivalences of the quotient of a minimal smooth prorepresentable groupoid are isomorphisms}
\author{The Stacks Project Authors}
\date{stacks-project:formal-defos.tex:local-snapshot-20261012; GNU FDL 1.2 or later}
\begin{document}
\maketitle
\noindent\textit{Editorial note.} All delivered mathematical text is the author's own, copied from the pinned Stacks Project chapter file formal-defos.tex; no statement or proof step was rewritten, reordered or removed.\par\medskip
\noindent\textit{Editorial note.} The counted proof is the author's complete argument: the translation of the self-equivalence into a morphism of groupoids in functors, the tangent space square combined with minimality making the first component an isomorphism, and the treatment of the second component through the exact sequence or the fibre product description.\par\medskip
\noindent\textit{Editorial note.} The label lemma-minimal-prorepresentable-groupoid-autoequivalence names the counted statement, and the proof environment immediately following it in the source is delivered verbatim as its proof.\par\medskip
\noindent\textit{Editorial note.} The counted statement is carried by the author's own lemma environment in the chapter source.\par\medskip
\noindent\textit{Editorial note.} Pre-open gate: python operations/agent-authoring/tools/precard\_gate.py --source formal-defos.tex --label lemma-minimal-prorepresentable-groupoid-autoequivalence --start 7140 --end 7177 --local-id 28028 returned RESULT=OK (L28028 free, occupancy FREE, public interval clear). Candidate taken from sources/stacks-formal-defos/candidates-halfA.json.\par\medskip
\noindent\textit{Editorial note.} The statement and proof are copied from the pinned chapter file; the cited inputs are kept as external pointers with their canonical Stacks tags.\par\medskip
\noindent\textit{Allowed standard prerequisites (editorial).}\begin{itemize}
\item Morphisms of the quotient are morphisms of groupoids in functors. Same chapter, the definition of a groupoid in functors and the equivalence between morphisms of the quotient and morphisms of the groupoid. Supplies the translation of a self-equivalence into a pair of morphisms phi : U -> U and psi : R -> R.
\item Characterization of minimality through tangent spaces. Same chapter, the criterion describing minimality of a smooth prorepresentable groupoid in terms of the tangent spaces and the action of the derivations of the residue field. Supplies the conclusion that the differential of phi is injective, hence bijective by dimension reasons.
\item Surjectivity criterion for morphisms of prorepresentable functors and the exact sequence computing the infinitesimal automorphisms. Same chapter, the lemma that a morphism of prorepresentable functors inducing an injection on tangent spaces is an isomorphism, and the exact sequence for the infinitesimal automorphism group of the quotient. Supply the conclusion that phi is an isomorphism and the alternative argument identifying psi with the induced morphism on the fibre product.
\end{itemize}\medskip
\ifdefined\hypertarget\hypertarget{sourcepointer-1}{}\fi
\noindent\textit{Source reference.} Original reference 06K3: sources/stacks-formal-defos/pinned-original-source.tex, source lines 6372--6372. Prerequisite pointer invoked at proof line(s) 7152; it names a standard fact of the same chapter that the counted argument uses, and every step of the counted proof built on it is retained verbatim in the delivered body.\par\medskip
\ifdefined\hypertarget\hypertarget{sourcepointer-2}{}\fi
\noindent\textit{Source reference.} Original reference 06KN: sources/stacks-formal-defos/pinned-original-source.tex, source lines 7079--7079. Prerequisite pointer invoked at proof line(s) 7164; it names a standard fact of the same chapter that the counted argument uses, and every step of the counted proof built on it is retained verbatim in the delivered body.\par\medskip
\ifdefined\hypertarget\hypertarget{sourcepointer-3}{}\fi
\noindent\textit{Source reference.} Original reference 06L6: sources/stacks-formal-defos/pinned-original-source.tex, source lines 6943--6943. Prerequisite pointer invoked at proof line(s) 7173; it names a standard fact of the same chapter that the counted argument uses, and every step of the counted proof built on it is retained verbatim in the delivered body.\par\medskip
\ifdefined\hypertarget\hypertarget{sourcepointer-4}{}\fi
\noindent\textit{Source reference.} Original reference 06KP: sources/stacks-formal-defos/pinned-original-source.tex, source lines 7116--7116. Prerequisite pointer invoked at proof line(s) 7167; it names a standard fact of the same chapter that the counted argument uses, and every step of the counted proof built on it is retained verbatim in the delivered body.\par\medskip
\begin{lemma}
\label{lemma-minimal-prorepresentable-groupoid-autoequivalence}
Let $(U, R, s, t, c)$ be a minimal smooth prorepresentable groupoid in
functors on $\mathcal{C}_\Lambda$. If $\varphi : [U/R] \to [U/R]$ is an
equivalence of categories cofibered in groupoids, then $\varphi$ is an
isomorphism.
\end{lemma}

\begin{proof}
A morphism $\varphi : [U/R] \to [U/R]$ is the same thing as a
morphism $\varphi : (U, R, s, t, c) \to (U, R, s, t, c)$ of
groupoids in functors over $\mathcal{C}_\Lambda$ as defined in
Definition \ref{definition-groupoid-in-functors}.
Denote $\phi : U \to U$ and $\psi : R \to R$ the corresponding morphisms.
Because the diagram
$$
\xymatrix{
& \text{Der}_\Lambda(k, k) \ar[dr]_\gamma \ar[dl]^\gamma \\
TU \ar[rr]_{d\phi} \ar[d] & & TU \ar[d]  \\
T[U/R] \ar[rr]^{d\varphi} & & T[U/R]
}
$$
is commutative, since $d\varphi$ is bijective, and since we have
the characterization of minimality in
Lemma \ref{lemma-characterize-minimal-groupoid-in-functors}
we conclude that $d\phi$ is injective (hence bijective by dimension reasons).
Thus $\phi : U \to U$ is an isomorphism by
Lemma \ref{lemma-surjective-morphism-prorepresentable-functor}.
We can use a similar argument, using the exact sequence
$$
0 \to \text{Inf}([U/R]) \to TR \to TU \oplus TU
$$
of
Lemma \ref{lemma-deformation-functor-diagonal}
to prove that $\psi : R \to R$ is an isomorphism. But is also a consequence
of the fact that $R = U \times_{[U/R]} U$ and that $\varphi$ and $\phi$
are isomorphisms.
\end{proof}

\end{document}
